% !TeX program = xelatex
% Απαιτεί XeLaTeX ή LuaLaTeX (λόγω fontspec/polyglossia) — ΔΕΝ μεταγλωττίζεται με pdflatex.
% Η γραμματοσειρά "Linux Libertine O"/"Linux Biolinum O" περιλαμβάνεται σε κάθε πλήρη
% εγκατάσταση TeX Live/MiKTeX και στο Overleaf. Αν ανεβάσετε το αρχείο στο Overleaf,
% ανεβάστε μαζί και τις εικόνες από results/ (βλ. \graphicspath παρακάτω) ή αλλάξτε το
% path ώστε να δείχνει εκεί που τις έχετε.
\documentclass[11pt,a4paper]{article}

\usepackage{fontspec}
\usepackage{polyglossia}
\setmainlanguage{greek}
\setotherlanguage{english}
\setmainfont{Linux Libertine O}
\setsansfont{Linux Biolinum O}
\setmonofont{DejaVu Sans Mono}[Scale=0.9]

\usepackage[a4paper,margin=2.6cm]{geometry}
\usepackage{parskip}
\usepackage{amsmath,amssymb}
\usepackage{booktabs}
\usepackage{longtable}
\usepackage{array}
\usepackage{graphicx}
\graphicspath{{../results/}}
\usepackage{enumitem}
\usepackage{caption}
\usepackage[hidelinks,colorlinks=true,linkcolor=blue!50!black,citecolor=blue!50!black,urlcolor=blue!50!black]{hyperref}
\usepackage{xcolor}

\newcommand{\file}[1]{\texttt{#1}}
\newcommand{\var}[1]{\texttt{#1}}

% Επιτρέπει αλλαγή γραμμής μετά από κάθε underscore, ώστε τα μακριά ονόματα αρχείων
% να μη βγαίνουν έξω από το περιθώριο ή τη στήλη ενός πίνακα.
\let\oldunderscore\_
\renewcommand{\_}{\oldunderscore\allowbreak}
\setlength{\emergencystretch}{3em}

\setlist[itemize]{topsep=2pt,itemsep=2pt}
\setlist[enumerate]{topsep=2pt,itemsep=2pt}

\title{\textbf{Online Convex Optimization based tuning of\\ Mixture of Experts for Sports Outcome Prediction}\\[6pt]
\large Αναφορά προόδου: θεωρία, υλοποίηση και ευρήματα}
\author{[Ονοματεπώνυμο φοιτητή]}
\date{\today}

\begin{document}
\maketitle
\tableofcontents
\newpage

% =====================================================================
\section{Εισαγωγή}
% =====================================================================

Η διπλωματική αυτή αντιμετωπίζει το πρόβλημα πρόβλεψης αποτελέσματος ποδοσφαιρικών
αγώνων (Νίκη Γηπεδούχου / Ισοπαλία / Νίκη Φιλοξενούμενου, «1X2») όχι ως πρόβλημα
εκπαίδευσης ενός μοντέλου πάνω σε στατιστικά αγώνων -- κάτι που τα ίδια τα στοιχηματικά
sites κάνουν ήδη, με πολύ περισσότερα δεδομένα και υπολογιστική ισχύ από όσα διαθέτει μια
διπλωματική -- αλλά ως πρόβλημα \emph{online μάθησης με συμβουλή ειδικών} (\emph{prediction
with expert advice}). Κάθε bookmaker αντιμετωπίζεται σαν έναν «expert» που δίνει, πριν από
κάθε αγώνα, τις δικές του πιθανότητες $(p_H, p_D, p_A)$ για τα τρία ενδεχόμενα. Ο στόχος
είναι να μάθουμε \emph{online}, αγώνα-αγώνα, μέσω αλγορίθμων \emph{Online Convex
Optimization} (OCO) -- Hedge, Online Gradient Descent (OGD), Follow-The-Regularized-Leader
(FTRL) -- ένα διάνυσμα βαρών πάνω στους experts, έτσι ώστε το σταθμισμένο μείγμα να
ξεπερνά κάθε μεμονωμένο bookmaker, αξιολογούμενο έναντι των ίδιων των αγορών (όχι έναντι
κάποιου «ground truth» μοντέλου).

Αυτή η αναφορά:
\begin{itemize}
  \item παρουσιάζει το θεωρητικό υπόβαθρο κάθε αλγορίθμου/τεχνικής που χρησιμοποιείται,
  \item περιγράφει τι έχει χτιστεί μέχρι στιγμής και γιατί, με τη σειρά που χτίστηκε,
  \item συνοψίζει τα ευρήματα κάθε ενότητας,
  \item και περιλαμβάνει πλήρη «χάρτη» του κώδικα: ποιο αρχείο κάνει τι, τι εισάγει από
        ποιο άλλο αρχείο, και τι αρχεία αποτελεσμάτων παράγει (Παράρτημα~\ref{app:filemap}).
\end{itemize}

Ο κώδικας είναι μια ακολουθία από ανεξάρτητα Python scripts στον φάκελο \file{scripts/},
το καθένα διαβάζει/γράφει επίπεδα αρχεία CSV/PNG στους φακέλους \file{data/} και
\file{results/}. Δεν υπάρχει build step, linter ή test suite -- πρόκειται για ένα project
ανάλυσης δεδομένων/ερευνητικών πειραμάτων, όχι συμβατικό λογισμικό.

% =====================================================================
\section{Δεδομένα}
% =====================================================================

\subsection{Πηγή και εύρος}

Τα δεδομένα προέρχονται από το \url{football-data.co.uk}, ενότητα «Main Leagues» (η
ενότητα «Extra Leagues» έχει διαφορετική μορφή και συνήθως καμία στήλη αποδόσεων, άρα
αποκλείεται). Καλύπτονται:
\begin{itemize}
  \item \textbf{22 πρωταθλήματα}: Αγγλία (5 κατηγορίες), Σκωτία (4), Γερμανία/Ιταλία/
        Ισπανία/Γαλλία (2 η καθεμία), Ολλανδία, Βέλγιο, Πορτογαλία, Τουρκία, Ελλάδα
        (1 η καθεμία),
  \item \textbf{10 σεζόν}: 2016/17 έως 2025/26,
  \item \textbf{76.639 αγώνες}, από τους οποίους οι 76.584 έχουν τουλάχιστον μία χρήσιμη
        απόδοση bookmaker -- αυτοί είναι οι γύροι πάνω στους οποίους τρέχουν οι αλγόριθμοι,
  \item \textbf{26 bookmakers/«experts»} μετά τον καθαρισμό (βλ.~\S\ref{sec:dataproc}).
\end{itemize}

Το project ξεκίνησε σε μικρότερη κλίμακα (5 σεζόν, 10 πρωταθλήματα) και επεκτάθηκε στην
πορεία· αρκετά ευρήματα άλλαξαν ανάμεσα στις δύο κλίμακες (βλ.~\S\ref{sec:pooling}).

\subsection{Επεξεργασία δεδομένων}
\label{sec:dataproc}

Το \file{data\_processer.py} διαβάζει κάθε ωμό (raw) αρχείο CSV, εντοπίζει κάθε τριάδα
στηλών αποδόσεων \{prefix\}H/\{prefix\}D/\{prefix\}A έναντι μιας hardcoded λίστας γνωστών
bookmakers, και μετατρέπει τις ωμές δεκαδικές αποδόσεις σε καθαρές
πιθανότητες αφαιρώντας το περιθώριο του bookmaker (\emph{de-vig}):
\begin{equation}
  p_i = \frac{1/O_i}{\sum_{j\in\{H,D,A\}} 1/O_j}, \qquad
  \text{overround} = \sum_{j} \frac{1}{O_j} - 1,
\end{equation}
όπου $O_i$ η δεκαδική απόδοση του ενδεχομένου $i$. Το overround (περιθώριο/vig) είναι το
πόσο πάνω από το «δίκαιο» 100\% αθροίζουν οι ωμές πιθανότητες -- το κέρδος του bookmaker.

Δύο ιδιαιτερότητες των δεδομένων αντιμετωπίζονται ρητά:
\begin{itemize}
  \item \textbf{Rebrand VC$\to$BetVictor}: το «VC Bet» μετονομάστηκε σε «BetVictor» (BV) στη
        μέση της χρονοσειράς (VC: σεζόν 2021/22--2023/24, BV: από 2025/26), χωρίς ποτέ να
        συνυπάρχουν στην ίδια σεζόν -- συγχωνεύονται σε μία συνεχή ταυτότητα
        \var{VC\_BV} (opening) / \var{VC\_BVC} (closing).
  \item \textbf{Στήλες συγκεντρωτικών αποδόσεων} (\var{Max}, \var{Avg}, και το παλαιότερο
        \var{BbAv}/\var{BbMx} του Betbrain, ορατό μόνο σε σεζόν πριν το $\sim$2021)
        αποκλείονται από το panel των experts -- είναι ήδη συνδυασμοί άλλων bookmakers,
        όχι ανεξάρτητοι experts.
\end{itemize}

\subsection{Opening vs.\ Closing αποδόσεις}
\label{sec:openclose}

Το football-data.co.uk δίνει, για τους περισσότερους bookmakers, \textbf{δύο} τιμές: μια
\emph{opening} (στήλες \{book\}H/D/A, π.χ.\ \var{B365H/D/A}) τη στιγμή που ανοίγει η
αγορά, και μια \emph{closing} (στήλες \{book\}CH/CD/CA, π.χ.\ \var{B365CH/CD/CA}) λίγο πριν
το ματς, αφού έχει ήδη κινηθεί όλο το «έξυπνο χρήμα». Το πρόθεμα ενός closing bookmaker
λήγει πάντα σε «C» -- σύμβαση που διατηρείται αυτούσια από το \file{data\_processer.py}, ακόμη
και μέσα από το merge VC/BV (\var{VC\_BV} = opening, \var{VC\_BVC} = closing). Το τρέχον
panel των 26 experts χωρίζεται έτσι \textbf{καθαρά} σε 13 opening / 13 closing, ένα ζευγάρι
ανά υποκείμενο bookmaker, χωρίς αταίριαστα ονόματα (βλ.~\S\ref{sec:openclose-results}).

% =====================================================================
\section{Θεωρητικό υπόβαθρο}
% =====================================================================

\subsection{Online Convex Optimization και πρόβλεψη με συμβουλή ειδικών}

Στο πλαίσιο \emph{Online Convex Optimization} (OCO), σε κάθε γύρο $t=1,\dots,T$ ένας
learner επιλέγει μια απόφαση $w_t$ μέσα σε ένα κυρτό σύνολο $\mathcal{K}$, στη συνέχεια
αποκαλύπτεται μια κυρτή συνάρτηση απώλειας $\ell_t(\cdot)$ και ο learner υφίσταται ζημία
$\ell_t(w_t)$. Ο στόχος δεν είναι να ελαχιστοποιηθεί η απόλυτη ζημία, αλλά η
\textbf{μεταμέλεια} (\emph{regret}) έναντι της καλύτερης σταθερής επιλογής εκ των υστέρων:
\begin{equation}
  R_T = \sum_{t=1}^T \ell_t(w_t) - \min_{w\in\mathcal{K}} \sum_{t=1}^T \ell_t(w).
\end{equation}
Στην ειδική περίπτωση «πρόβλεψης με συμβουλή ειδικών» (\emph{prediction with expert
advice}), το $\mathcal{K}$ είναι ο simplex πιθανοτήτων πάνω από $N$ experts, και
$\ell_t(w) = \sum_i w_i \, \ell_{i,t}$ όπου $\ell_{i,t}$ η ζημία (εδώ: log-loss) του expert
$i$ στον γύρο $t$. Κάθε αλγόριθμος παρακάτω (Hedge/OGD/FTRL) είναι μια διαφορετική
στρατηγική επιλογής $w_t$ με εγγυημένο άνω φράγμα στο $R_T$ (τυπικά $O(\sqrt{T})$),
ανεξάρτητα από το πώς παράγεται η ακολουθία $\ell_t$ -- ακόμη κι αν είναι αντίπαλη
(\emph{adversarial}). Αυτό είναι και το βασικό πλεονέκτημα έναντι ενός στατικού μοντέλου: η
αγορά στοιχημάτων δεν είναι στάσιμη (βλ.~\S\ref{sec:temporal}), και οι αλγόριθμοι OCO δεν
υποθέτουν στασιμότητα.

Σε κάθε γύρο, η πρόβλεψη μείγματος είναι $\hat p_t = \sum_i w_{i,t}\, p_{i,t}$ όπου
$p_{i,t}\in\Delta^2$ οι πιθανότητες που έδωσε ο expert $i$, και η ζημία μετριέται με
log-loss: $\ell_{i,t} = -\log p_{i,t}(y_t)$, όπου $y_t\in\{H,D,A\}$ το πραγματικό
αποτέλεσμα.

\subsection{Hedge (Multiplicative Weights)}

Ο Hedge διατηρεί ένα θετικό βάρος $w_i$ ανά expert και το πολλαπλασιάζει εκθετικά με βάση
τη ζημία:
\begin{equation}
  w_i(t+1) = w_i(t)\, \exp\!\big(-\eta_t\, \ell_{i,t}\big), \qquad
  \eta_t = \frac{\sqrt{\ln N / (t+1)}}{M},
\end{equation}
όπου $M=\ln(1/\varepsilon)$ ένα άνω φράγμα στη ζημία ενός γύρου (με $\varepsilon$ ένα κάτω
όριο αποκοπής στις πιθανότητες, ώστε το log-loss να μη γίνεται άπειρο). Η πρόβλεψη
προκύπτει κανονικοποιώντας: $v_i(t) = w_i(t)/\sum_j w_j(t)$.

\subsection{Online Gradient Descent (OGD)}

Ο OGD κάνει ένα βήμα κλίσης και προβάλλει πίσω στον simplex $\Delta$:
\begin{equation}
  w(t+1) = \Pi_{\Delta}\big(w(t) - \eta_t\, \nabla \ell_t(w(t))\big), \qquad
  \eta_t = \frac{D}{M\sqrt{t+1}},
\end{equation}
με $D=\sqrt{2}$ (διάμετρος του simplex) και $\Pi_\Delta$ την Ευκλείδεια προβολή στον simplex
(αλγόριθμος Duchi et al., 2008). Είναι «greedy»: το επόμενο βήμα ξεκινά από το \emph{τρέχον}
$w(t)$.

\subsection{Follow-The-Regularized-Leader (FTRL)}

Ο FTRL είναι «lazy»: σε κάθε γύρο ξαναϋπολογίζει τα βάρη από το \textbf{σύνολο} της
αθροιστικής ζημίας μέχρι τώρα, όχι από το προηγούμενο βήμα:
\begin{equation}
  w(t) = \Pi_{\Delta}\Big(w_0 - \eta_t \sum_{s\le t} \ell_s\Big), \qquad
  \eta_t = \frac{D}{M\sqrt{t+1}}.
\end{equation}

\subsection{Η αναγωγή Sleeping Experts / Specialists}
\label{sec:sleeping}

Δεν προσφέρει κάθε bookmaker απόδοση για κάθε αγώνα -- πολλοί καλύπτουν μόνο ένα υποσύνολο
αγώνων/σεζόν (\emph{partial coverage}). Οι κλασικοί Hedge/OGD/FTRL υποθέτουν ότι όλοι οι
experts μιλούν σε κάθε γύρο, οπότε γενικεύονται μέσω της αναγωγής \textbf{Sleeping
Experts / Specialists} (Freund, Schapire, Singer \& Warmuth, 1997· Blum \& Mansour, 2007):
\begin{itemize}
  \item κάθε expert κρατά τη δική του μόνιμη κατάσταση (βάρος για Hedge, θέση για OGD,
        αθροιστική ζημία για FTRL)·
  \item σε κάθε γύρο $t$ μόνο το «ξύπνιο» υποσύνολο $A(t)$ (οι bookmakers που έδωσαν
        απόδοση σε αυτόν τον αγώνα) χρησιμοποιείται, με επανακανονικοποίηση \emph{μόνο}
        πάνω στο $A(t)$·
  \item ενημερώνεται η κατάσταση \emph{μόνο} των ξύπνιων experts· η κατάσταση των
        «κοιμισμένων» παγώνει ακριβώς όπως ήταν.
\end{itemize}
Αυτό δίνει σε κάθε expert $i$ εγγύηση μεταμέλειας μετρημένη \emph{μόνο} πάνω στους γύρους
που ήταν πραγματικά ξύπνιος, χωρίς να απαιτείται συμμετοχή σε κάθε γύρο. Υλοποιείται στο
\file{sleeping\_experts.py} και είναι το θεμέλιο πάνω στο οποίο χτίζεται σχεδόν το
υπόλοιπο project.

\subsection{Fixed-Share}

Παρατηρήθηκε ότι σπάνιοι (π.χ.\ μονής σεζόν) bookmakers μπορούν να καταλήξουν με
δυσανάλογα μεγάλο τελικό βάρος Hedge, επειδή τίποτα δεν διαβρώνει το βάρος τους όσο
κοιμούνται. Το Fixed-Share (Herbster \& Warmuth, 1998) προσθέτει μια μικρή, σταθερή διαρροή
βάρους προς την ομοιόμορφη κατανομή σε κάθε γύρο που ο expert είναι ξύπνιος:
\begin{equation}
  w_i(t+1) = (1-\alpha)\, w_i^{\text{Hedge}}(t+1) \;+\; \alpha \cdot \frac{1}{|A(t)|}.
\end{equation}
Για $\alpha=0$ ανάγεται ακριβώς στον απλό sleeping-experts Hedge (επαληθεύεται με assertion
στο \file{fixed\_share.py}).

\subsection{Βαθμονόμηση: online temperature scaling ως δεύτερο στάδιο OCO}
\label{sec:calib-theory}

Η ανάλυση βαθμονόμησης (\file{calibration\_analysis.py}) αποκάλυψε μεροληψία τύπου
\emph{favorite-longshot bias}, κοινή σε όλες τις σειρές (και τα μείγματα): οι χαμηλές
προβλεπόμενες πιθανότητες είναι συστηματικά υψηλότερες από την πραγματική συχνότητα, και οι
υψηλές συστηματικά χαμηλότερες -- δηλαδή \emph{underconfidence}. Η διόρθωση είναι
\emph{temperature scaling} με εκθέτη $s=1/T$:
\begin{equation}
  \hat p_k^{(s)} = \frac{\hat p_k^{\,s}}{\sum_j \hat p_j^{\,s}}, \qquad s>1 \Rightarrow
  \text{«οξύνει» τις πιθανότητες μακριά από το ομοιόμορφο.}
\end{equation}
Με $\log\hat p = (\log\hat p_1,\log\hat p_2,\log\hat p_3)$, η ζημία
\begin{equation}
  \ell(s) = -s\,\log\hat p_y + \log\!\Big(\sum_k \exp(s\,\log\hat p_k)\Big)
\end{equation}
είναι ένα log-sum-exp μιας γραμμικής συνάρτησης του $s$, άρα \textbf{κυρτή} στο $s$ -- οπότε
μπορεί να μαθευτεί online με απλό προβαλλόμενο OGD πάνω στο διάστημα $[S_{\min}, S_{\max}]$,
χωρίς νέο μηχανισμό: είναι η ίδια συνταγή OCO, απλά με μεταβλητή απόφασης 1 διάστασης αντί
για simplex. Αυτό είναι ένα \textbf{δεύτερο στάδιο OCO}, πάνω στην ήδη-μαθημένη πρόβλεψη
μείγματος του πρώτου σταδίου.

Το πρωτόκολλο εκπαίδευσης είναι πάντα χρονολογικό, χωρίς διαρροή πληροφορίας: ο ρυθμός
μάθησης επιλέγεται μόνο στο πρώτο \var{TRAIN\_FRACTION} του ιστορικού μιας σειράς, και η
τελική αξιολόγηση γίνεται είτε μόνο στο held-out val κομμάτι
(\file{calibration\_correction.py}) είτε, για δίκαιη σύγκριση raw έναντι calibrated, σε ΕΝΑ
συνεχές πέρασμα πάνω σε ολόκληρο το ιστορικό με τον ήδη επιλεγμένο ρυθμό
(\file{final\_ranking.py}).

\subsubsection{Per-season και warm-start παραλλαγές}

Το global fit παραπάνω χρησιμοποιεί σχήμα ρυθμού μάθησης $\eta_t = c/\sqrt{t+1}$ που
συρρικνώνεται σε όλη τη 10ετία -- μέχρι τις τελευταίες σεζόν το $s$ είναι ουσιαστικά
«παγωμένο» γύρω από έναν συμβιβασμό μακράς διάρκειας, χωρίς να προλαβαίνει μια μεταβολή
που αλλάζει σεζόν-σεζόν (βλ.~\S\ref{sec:temporal}). Δύο εναλλακτικές εξετάστηκαν:
\begin{itemize}
  \item \textbf{Per-season} (\file{calibration\_per\_season.py}): μηδενισμός \emph{και} του
        μετρητή γύρου $t$ \emph{και} του $s$ (πίσω σε $s=1$) στην αρχή κάθε σεζόν --
        γρήγορη προσαρμογή, αλλά με κόστος «κρύου ξεκινήματος» 10 φορές, πάνω σε
        μικρότερο δείγμα ανά σεζόν ($\sim$7.658 γύροι έναντι 76.584).
  \item \textbf{Warm-start} (\file{calibration\_warmstart.py}): μηδενισμός \emph{μόνο} του
        μετρητή γύρου (ώστε το $\eta_t$ να παραμένει γρήγορο), αλλά το $s$ κληρονομείται
        από το τελικό fitted $s$ της προηγούμενης σεζόν αντί να μηδενίζεται.
\end{itemize}

\subsection{Στατιστική σημαντικότητα: Moving Block Bootstrap}
\label{sec:bootstrap-theory}

Οι διαδοχικές ανά-αγώνα ζημίες \textbf{δεν} είναι ανεξάρτητα δείγματα (ένα σαββατοκύριακο
με πολλές ανατροπές χτυπά τη ζημία όλων των forecasters ταυτόχρονα), άρα ένα απλό
$t$-test ή i.i.d.\ bootstrap θα υποεκτιμούσε τη διακύμανση. Το \textbf{moving block
bootstrap} (Künsch, 1989) δειγματοληπτεί \textbf{συνεχόμενα μπλοκ} γύρων αντί μεμονωμένων
γύρων, διατηρώντας τη βραχυπρόθεσμη συσχέτιση. Το μήκος μπλοκ ακολουθεί τον εμπειρικό
κανόνα $n^{1/3}$ (Hall, Horowitz \& Jing, 1995). Για κάθε σύγκριση, ορίζεται
$\text{diff}_t = \ell_t(A) - \ell_t(B)$ πάνω στους κοινούς έγκυρους γύρους των δύο σειρών,
σε χρονολογική σειρά· ένα \textbf{αρνητικό} μέσο diff σημαίνει ότι η πρώτη σειρά $A$ είναι
\textbf{καλύτερη} (χαμηλότερο log-loss). Αυτή είναι εννοιολογικά η ίδια ιδέα με το κλασικό
τεστ Diebold--Mariano, μέσω bootstrap αντί ασυμπτωτικής κανονικής προσέγγισης.

\subsection{Mann-Kendall trend test}

Για τον έλεγχο τάσης σε σειρές μήκους μόλις 10 σεζόν, χρησιμοποιείται το μη-παραμετρικό
τεστ Mann--Kendall αντί προσαρμογής ευθείας OLS: μετρά συμφωνούντα έναντι ασύμφωνων ζευγών
σε όλη τη σειρά,
\begin{equation}
  S=\sum_{k<j}\mathrm{sign}(x_j-x_k), \qquad
  Z = \frac{S \mp 1}{\sqrt{\operatorname{Var}(S)}},
\end{equation}
χωρίς να υποθέτει γραμμικότητα ή κανονικά κατάλοιπα -- καλύτερη επιλογή από OLS $p$-values
για τόσο κοντή σειρά.

\subsection{Κριτήριο Kelly (value betting)}
\label{sec:kelly-theory}

Για δεδομένη εκτιμώμενη πιθανότητα $p$ και δεκαδική απόδοση $O$, το κλάσμα Kelly (ποσοστό
τράπεζας που στοιχηματίζεται) είναι
\begin{equation}
  f^{*} = \frac{p\,O - 1}{O - 1},
\end{equation}
το κλάσμα που μεγιστοποιεί τον αναμενόμενο λογαριθμικό ρυθμό ανάπτυξης της τράπεζας
μακροπρόθεσμα. Στο \file{value\_betting.py} χρησιμοποιείται πολλαπλάσιο του πλήρους Kelly
(\var{kelly\_mult}) πάνω στο ενδεχόμενο με το υψηλότερο θετικό expected value ανά γύρο.

% =====================================================================
\section{Αρχιτεκτονική / σειρά υλοποίησης}
\label{sec:architecture}
% =====================================================================

Όλα τα downstream scripts διαβάζουν το \file{data/processed/odds\_long.csv} (χτισμένο από
το \file{data\_processer.py}). Σχεδόν κάθε script μετά από αυτό εισάγει απευθείας κοινό κώδικα
από το \file{sleeping\_experts.py}. Η δομή εξαρτήσεων:

{\small
\begin{verbatim}
download_data.ps1
   |
   v
data_processer.py  <-- παράγει ΚΑΙ τα δύο datasets σε ένα run
   |
   +--> data/processed/odds_long.csv
   |         |
   |         v
   |    sleeping_experts.py  <-- μόνο algorithms/data-loading, χωρίς main();
   |         |                   όλα τα παρακάτω εισάγουν από αυτό
   |         +--> sleeping_experts_experiment.py  (παράγει results_table.csv)
   |         |
   |         +--> calibration_analysis.py
   |         |         |
   |         |         v
   |         +--> calibration_correction.py
   |         |         |
   |         |         v
   |         +--> final_ranking.py  <-- calibrate_full_history() εισάγεται
   |         |         |                από 6 αρχεία:
   |         |         +--> significance_test.py
   |         |         +--> value_betting.py
   |         |         |         +--> value_betting_online_kelly.py
   |         |         +--> market_devig_comparisons.py
   |         |         +--> contextual_experts.py
   |         |         +--> pooling_window_comparison.py
   |         |         +--> calibration_per_season.py
   |         |                   |
   |         |                   v
   |         |              calibration_warmstart.py
   |         |
   |         +--> fixed_share.py
   |         +--> outcome_class_accuracy.py
   |         +--> temporal_trends.py
   |
   +--> data/processed/odds_long_shin.csv  (διαβάζεται απευθείας από τα
                                             value_betting.py και
                                             market_devig_comparisons.py)
\end{verbatim}
}

Το παλιότερο δίπτυχο \file{process\_data.py} + \file{process\_data\_shin.py} (χωριστό
script ανά normalization) αντικαταστάθηκε από το ενιαίο \file{data\_processer.py}·
ομοίως, 8 παλιότερα, χωριστά scripts σύγκρισης/value-betting (\file{opening\_vs\_closing.py},
\file{shin\_vs\_basic\_comparison.py}, \file{shin\_vs\_basic\_calibration.py},
\file{shin\_vs\_basic\_value\_betting.py}, \file{opening\_closing\_quality\_transfer.py},
\file{value\_betting\_all\_bookmakers.py}, \file{value\_betting\_closing\_trained.py},
\file{value\_betting\_training\_panel\_sweep.py}) ενοποιήθηκαν στα \file{value\_betting.py}
και \file{market\_devig\_comparisons.py} -- επαληθεύτηκε τιμή-για-τιμή ταυτόσημη έξοδος πριν
την αλλαγή σε κάθε περίπτωση· τα παλιά αρχεία παραμένουν στο \file{scripts/deprecated/} μόνο
για αναφορά.

Κάθε script τρέχει ανεξάρτητα (\texttt{cd scripts; python <name>.py}), και τίποτα δεν είναι
cache-αρισμένο ανάμεσα σε scripts εκτός από τα CSV που το καθένα γράφει στο \file{results/}
και κάποιο επόμενο διαβάζει πίσω -- κάθε script ξανα-υπολογίζει τους δικούς του αλγορίθμους
από την αρχή σε κάθε εκτέλεση.

% =====================================================================
\section{Τι έχει χτιστεί και τι βρήκαμε}
% =====================================================================

\subsection{Πυρήνας αλγορίθμων: Hedge / OGD / FTRL έναντι bookmakers}
\label{sec:core-algos}

Η πρώτη υλοποίηση έγινε σε ένα στενό panel 6 bookmakers πλήρους κάλυψης (πρώιμο draft που
δεν περιλαμβάνεται σε αυτό το project), και στη συνέχεια γενικεύτηκε σε ολόκληρο το panel
των 26 experts με μερική κάλυψη μέσω της αναγωγής Sleeping Experts
(\file{sleeping\_experts.py} -- \file{results\_table.csv}). Στο πλήρες σύνολο δεδομένων
(76.584 γύροι):
\begin{itemize}
  \item Το \textbf{OGD είναι ο καλύτερος} από τους 3 αλγορίθμους πυρήνα. Και οι τρεις
        ξεπερνούν σημαντικά την απλή ομοιόμορφη μέση τιμή (uniform average) στο πλήρες
        10ετές dataset· με μισό χρονικό παράθυρο το πετυχαίνει μόνο το OGD
        (βλ.~\S\ref{sec:pooling} για την απομόνωση αυτής της επίδρασης βάθους δεδομένων).
  \item Το \textbf{PSC (Pinnacle closing) έχει στατιστικά σημαντικό προβάδισμα έναντι του
        OGD σε αυτό το panel}, πριν και μετά τη βαθμονόμηση. Στη διαμόρφωση αναφοράς
        (\S\ref{sec:bestconfig}) το προβάδισμα αυτό παύει να είναι σημαντικό.
  \item Log-loss στο πλήρες panel (26 experts): OGD 0{,}99745, Hedge 0{,}99845, FTRL
        0{,}99850, uniform average 0{,}99886 -- όλα χειρότερα από το PSC (0{,}99714), τον
        καλύτερο μεμονωμένο bookmaker.
\end{itemize}

\begin{figure}[htbp]
  \centering
  \includegraphics[width=0.92\linewidth]{sleeping_experts_comparison.png}
  \caption{Σύγκριση αλγορίθμων πυρήνα έναντι του καλύτερου bookmaker πλήρους κάλυψης --
  τάση στον χρόνο (αριστερά) και τελική κατάταξη (δεξιά).}
\end{figure}

\subsection{Βαθμονόμηση}

Η ανάλυση βαθμονόμησης (\file{calibration\_analysis.py}, Σχήμα~\ref{fig:calib}) διέγνωσε το
favorite-longshot bias περιγραφόμενο στο \S\ref{sec:calib-theory}. Η διόρθωση μέσω online
temperature scaling (\file{calibration\_correction.py}, \file{final\_ranking.py}) βοηθά
\textbf{σημαντικά και συστηματικά}, σε κάθε αλγόριθμο και σε κάθε χρονικό παράθυρο που
δοκιμάστηκε.

\begin{figure}[htbp]
  \centering
  \includegraphics[width=0.92\linewidth]{calibration_reliability.png}
  \caption{Διάγραμμα αξιοπιστίας (reliability diagram) και «κενό» βαθμονόμησης, πριν τη
  διόρθωση.}
  \label{fig:calib}
\end{figure}

\paragraph{Per-season και warm-start.} Το per-season fit (\S\ref{sec:calib-theory}) είναι
σημαντικά καλύτερο από το raw, αλλά στατιστικά αδιάκριτο από το ήδη υπάρχον global fit --
το μικρότερο δείγμα ανά σεζόν πιθανότατα ακυρώνει το όφελος παρακολούθησης της
μετατόπισης (drift). Το warm-start variant που χτίστηκε σε αυτή τη συνεδρία εξαλείφει
οπτικά το «sawtooth» pattern του cold-start per-season fit (Σχήμα~\ref{fig:warmstart}), και
το log-loss του πέφτει \textbf{ανάμεσα} στο global και στο per-season για κάθε σειρά (π.χ.\
OGD: raw 0{,}99745, global 0{,}99693, per-season 0{,}99705, warm-start 0{,}99697) -- αλλά
είναι στατιστικά αδιάκριτο και από τα δύο (κανένα από τα «warm-start vs global» / «warm-start
vs per-season» δεν είναι σημαντικό σε κανένα από τα 11 σειρές που ελέγχθηκαν), ενώ ξεπερνά
σημαντικά το raw σχεδόν παντού. Δηλαδή είναι ένας γνήσιος συμβιβασμός, όχι αυστηρή
βελτίωση έναντι κάποιου από τα δύο υπάρχοντα variants -- η αφαίρεση του κόστους
«κρύου ξεκινήματος» δεν μεταφράστηκε σε μετρήσιμο κέρδος log-loss σε αυτό το μέγεθος
dataset.

\begin{figure}[htbp]
  \centering
  \includegraphics[width=0.8\linewidth]{calibration_warmstart.png}
  \caption{Fitted θερμοκρασία $T$ ανά σεζόν: cold-start reset (διακεκομμένη) έναντι
  warm-start (συνεχής) -- το warm-start εξομαλύνει το sawtooth pattern.}
  \label{fig:warmstart}
\end{figure}

\subsection{Στατιστική σημαντικότητα}

Το \file{significance\_test.py} τρέχει μια «μπαταρία» 12 συγκρίσεων μέσω moving block
bootstrap (\S\ref{sec:bootstrap-theory}). Στο τρέχον, πλήρες 10ετές dataset
(\file{significance\_test\_table.csv}), με αρνητική διαφορά να σημαίνει ότι το πρώτο-αναφερόμενο
είναι καλύτερο:
\begin{itemize}
  \item Και οι τρεις αλγόριθμοι ξεπερνούν σημαντικά την ομοιόμορφη μέση τιμή: OGD
        $-0{,}00141$, Hedge $-0{,}00041$, FTRL $-0{,}00037$ (όλα με $p<0{,}001$).
  \item Το OGD ξεπερνά σημαντικά τον Hedge ($-0{,}00100$), τον FTRL ($-0{,}00104$) και
        \emph{και} τον B365C ($-0{,}00063$), όλα με $p<0{,}001$.
  \item Το PSC είναι σημαντικά καλύτερο από το OGD \emph{σε αυτή τη διαμόρφωση} (πλήρες
        panel, proportional de-vig): $+0{,}00065$ πριν τη βαθμονόμηση
        (95\% CI $[0{,}00044,\,0{,}00086]$) και $+0{,}00043$ μετά. Στη διαμόρφωση αναφοράς
        (closing-only panel) η διαφορά παύει να είναι σημαντική --- βλ.\
        \S\ref{sec:bestconfig}.
  \item Η βαθμονόμηση βελτιώνει σημαντικά και τους 4 αλγορίθμους (raw μείον calibrated:
        $+0{,}00053$ έως $+0{,}00073$).
\end{itemize}
Με μισό χρονικό παράθυρο, το OGD είναι ο μόνος αλγόριθμος που ξεπερνά σημαντικά την
ομοιόμορφη μέση τιμή· το βάθος των 10 σεζόν είναι αυτό που φέρνει και τα άλλα δύο πάνω από
αυτήν (βλ.\ \S\ref{sec:pooling}).

\subsection{Ρύθμιση του βήματος μάθησης}
\label{sec:steptuning}

Το θεωρητικά βέλτιστο (worst-case) χρονοδιάγραμμα βήματος περιέχει το φράγμα
$M=\log(1/\varepsilon)=13{,}82$ για τη μέγιστη απώλεια ενός γύρου. Το φράγμα αυτό υποθέτει
ότι ένας bookmaker θα μπορούσε να τιμολογήσει το πραγματοποιηθέν αποτέλεσμα στο $10^{-6}$·
η μέγιστη απώλεια που όντως παρατηρείται στα δεδομένα είναι \textbf{3{,}76}. Το worst-case
βήμα είναι επομένως αρκετά πιο συντηρητικό απ' όσο δικαιολογούν τα δεδομένα, και αξίζει να
συντονιστεί εμπειρικά.

Το \file{learning\_rate\_study.py} σαρώνει έναν πολλαπλασιαστή βήματος και για τους τρεις
αλγορίθμους στο ίδιο πλέγμα, με επιλογή στο χρονολογικά πρώτο 75\% των γύρων και μέτρηση στο
υπόλοιπο 25\% --- το ίδιο πρωτόκολλο που χρησιμοποιείται ήδη για τον ρυθμό της βαθμονόμησης
(\S\ref{sec:calib-theory}). Αποτέλεσμα:

\begin{itemize}
  \item Τα \textbf{Hedge και FTRL είναι ήδη βέλτιστα} στο $\times 1$· κάθε αύξηση τα
        χειροτερεύει μονότονα.
  \item Το \textbf{OGD είναι βέλτιστο στο $\times 100$} (held-out $0{,}99765$ έναντι
        $0{,}99803$), με γνήσιο εσωτερικό βέλτιστο: τα $\times 200$, $\times 500$ και
        $\times 1000$ είναι χειρότερα, και train και held-out συμφωνούν στην επιλογή.
\end{itemize}

Η ασυμμετρία είναι δομική: στο OGD το $\eta$ πολλαπλασιάζει την κλίση \emph{ενός} γύρου, ενώ
στο FTRL πολλαπλασιάζει ολόκληρη τη συσσωρευμένη απώλεια, οπότε μεγάλο $\eta$ το καταρρέει
προς follow-the-leader. Επειδή και οι τρεις συντονίστηκαν με τον ίδιο τρόπο, η σύγκριση
μεταξύ τους παραμένει δίκαιη· ο συντελεστής του OGD είναι το
$\var{OGD\_STEP\_MULTIPLIER}$ στο \file{sleeping\_experts.py}.

Δοκιμάστηκαν και αρχές που δεν απαιτούν επιλογή παραμέτρου, με μέτρια αποτελέσματα: ένα
τίμιο a-priori φράγμα $M=\log(100)$ ανακτά μόλις το 15\% του κέρδους, το AdaGrad (scalar) το
20\%, ενώ το AdaGrad ανά expert είναι \emph{χειρότερο} από καθόλου ρύθμιση. Αξιοσημείωτο
είναι επίσης ότι η θεωρητικά ανώτερη εναλλακτική --- το log-loss είναι exp-concave, οπότε ο
Aggregating Algorithm (Vovk, 1990) πετυχαίνει regret $\le \ln N$ με σταθερό $\eta=1$ ---
αποδίδει στην πράξη \emph{χειρότερα από την ομοιόμορφη μέση τιμή} εδώ. Το φράγμα regret
αφορά στατικό comparator, ενώ το panel δεν είναι στατικό: η κάλυψη του PSC καταρρέει από
$\sim$100\% σε 38{,}8\% στην τελευταία σεζόν, που πέφτει μέσα στο held-out κομμάτι, και μια
posterior που συγκεντρώνεται γρήγορα τιμωρείται ακριβώς από έναν expert που εξαφανίζεται.

\subsection{Το μείγμα χωρίς την Pinnacle}
\label{sec:nopinnacle}

Το προφανές αντεπιχείρημα στο «το μείγμα ξεπερνά κάθε bookmaker εκτός του PSC» είναι ότι η
Pinnacle βρίσκεται ούτως ή άλλως \emph{μέσα} στο μείγμα, άρα το μείγμα ίσως απλώς την
αναπαράγει. Το \file{no\_pinnacle\_panel.py} το ελέγχει αφαιρώντας \textbf{και τις δύο}
στήλες Pinnacle (PS και PSC) από το panel και επανεκπαιδεύοντας από το μηδέν σε δύο panel:
το πλήρες (24 experts απομένουν, κάλυψη γύρων 100\%) και το closing-only (12 experts), που
είναι η διαμόρφωση αναφοράς του \S\ref{sec:bestconfig}.

Και στα δύο, κάθε σειρά βαθμολογείται στους \emph{ίδιους} γύρους --- την τομή όπου όλες οι
συγκρινόμενες σειρές ορίζονται --- με χρονικό διαχωρισμό 75/25 για μια ουρά που η ανάγνωση
των αποτελεσμάτων δεν έχει δει. Χωρίς αυτόν τον περιορισμό οι σειρές συγκρίνονται σε
διαφορετικά σύνολα αγώνων και η κατάταξη δεν σημαίνει τίποτα.

\paragraph{Πλήρες panel: το αποτέλεσμα αντέχει.} Σε 54.877 κοινούς γύρους, \textbf{και οι
τέσσερις αλγόριθμοι κατατάσσονται πάνω από όλους τους εναπομείναντες Tier-A bookmakers}:
Hedge $0{,}99772$ (βαθμονομημένο), FTRL $0{,}99778$, OGD $0{,}99782$, uniform $0{,}99837$,
έναντι VC\_BV $0{,}99987$ για τον καλύτερο bookmaker. Το OGD ξεπερνά \emph{σημαντικά} και
τους πέντε (περιθώρια $-0{,}00181$ έως $-0{,}00315$, όλα $p<0{,}001$), με πολύ φαρδύτερα
περιθώρια απ' ό,τι επιτυγχάνει έναντι του PSC στο πλήρες panel, και είναι πρώτο στην
held-out ουρά. Ο περιορισμός «ένας bookmaker μας ξεπερνά» αφορά επομένως \emph{μία}
εξαιρετική εταιρεία, όχι τη μέθοδο. Έχει και πρακτική σημασία: η Pinnacle δεν είναι
προσβάσιμη σε πολλές δικαιοδοσίες, οπότε για τον περισσότερο κόσμο η κρίσιμη σύγκριση είναι
ακριβώς αυτή --- και εκεί το μείγμα κερδίζει καθαρά.

\paragraph{Closing-only panel: εδώ η μάθηση εξαφανίζεται.} Στη διαμόρφωση αναφοράς η
αφαίρεση κοστίζει πολύ περισσότερο. Απομένουν 12 experts και μόνο 32.210 κοινοί γύροι, και
το αποτέλεσμα αλλάζει ποιοτικά:

\begin{itemize}
  \item Το OGD βγαίνει \#1 ($0{,}99919$ βαθμονομημένο) αλλά ξεπερνά σημαντικά μόνο τους
        \emph{τρεις} από τους πέντε (WHC $-0{,}00088$, IWC $-0{,}00077$, BWC $-0{,}00051$).
        Ο B365C ($p=0{,}505$) και ο VC\_BVC ($p=0{,}284$) είναι ισοπαλίες, και στην held-out
        ουρά ο VC\_BVC περνάει μπροστά από το OGD.
  \item Πιο διαφωτιστικό: το OGD ($0{,}99919$) γίνεται πρακτικά \emph{ταυτόσημο} με τον απλό
        μέσο όρο ($0{,}99921$), με τα Hedge ($0{,}99923$) και FTRL ($0{,}99932$) πίσω από
        αυτόν. Η μαθημένη στάθμιση παύει να προσφέρει οτιδήποτε.
\end{itemize}

Αυτό δεν αναιρεί το εύρημα του πλήρους panel, το εξειδικεύει --- και δένει με το κεντρικό
συμπέρασμα του \S\ref{sec:openclose-results}: η μάθηση αποδίδει όταν υπάρχει ποιοτική
διαφορά μεταξύ των experts. Το πλήρες panel την έχει (13 closing έναντι 13 συστηματικά
χειρότερων opening). Το closing-only panel τη χρωστά σε μεγάλο βαθμό στην παρουσία της
Pinnacle· μόλις φύγει, οι υπόλοιποι closing bookmakers είναι υπερβολικά όμοιοι μεταξύ τους
και δεν μένει τίποτα να μαθευτεί. Η καλύτερη διαμόρφωση της εργασίας, δηλαδή, είναι η
καλύτερη \emph{επειδή} περιέχει τον καλύτερο bookmaker της αγοράς --- τίμημα που δεν
κρύβεται, αφού η σύγκριση με τον PSC (\S\ref{sec:bestconfig}) γίνεται με τον PSC μέσα στο
μείγμα.

\subsection{Fixed-Share}

Το \file{fixed\_share.py} διορθώνει το πρόβλημα «παγωμένου βάρους»: σπάνιοι bookmakers
(συχνά μονής σεζόν) καταλήγουν με δυσανάλογα μεγάλο βάρος Hedge επειδή τίποτα δεν
το διαβρώνει όσο κοιμούνται. Μια μικρή διαρροή $\alpha=0{,}005$--$0{,}01$ περιορίζει το
πρόβλημα (π.χ.\ το βάρος του CLC την τελευταία φορά που ήταν ξύπνιος:
$0{,}240 \to 0{,}064$ για $\alpha=0{,}01$· οι τέσσερις bookmakers που ήταν ενεργοί μόνο στη
σεζόν 2024/25 --- 1XB, 1XBC, BF, BFC --- πέφτουν από $0{,}164$--$0{,}177$ σε ακριβώς
$0{,}100$, δηλαδή στην ισοτιμία $1/10$ του συνόλου των ξύπνιων εκείνου του γύρου), με μικρό
αλλά υπαρκτό κόστος στο συνολικό log-loss ($0{,}99845 \to 0{,}99885$, δηλαδή ελαφρώς
\emph{χειρότερο}).

Σημείωση μέτρησης (βλ.\ \S\ref{sec:bugfix}): το «βάρος» εδώ είναι το βάρος που κρατούσε ο
κάθε bookmaker \textbf{την τελευταία φορά που ήταν ξύπνιος}, όχι το βάρος στον τελευταίο
γύρο του dataset. Το δεύτερο είναι μηδέν για όποιον δεν έδωσε απόδοση στον τελευταίο αγώνα
(10 από τους 26 εδώ), και έκρυβε τα τέσσερα σοβαρότερα κρούσματα του φαινομένου.

\subsection{Contextual experts (ανά πρωτάθλημα)}
\label{sec:contextual}

Δοκιμάζει την αρχική ιδέα της διπλωματικής απευθείας: βοηθάει η εκμάθηση \textbf{χωριστών}
βαρών ανά πρωτάθλημα έναντι ενός ενιαίου (pooled) μοντέλου; Το \file{contextual\_experts.py}
βρίσκει ότι \textbf{όχι}. Πάνω σε 6 πρωταθλήματα πρώτης κατηγορίας (Premier League, Σκωτία,
Serie A, La Liga, Bundesliga, Ligue 1 -- σκόπιμα όχι και οι 22 κατηγορίες, ώστε να μη
συγκρίνονται ειδικοί σε πολύ διαφορετικά επίπεδα) το pooled μοντέλο ξεπερνά σημαντικά τους
ειδικούς συνολικά (μέση διαφορά $+0{,}00058$, 95\% CI $[+0{,}00023,\,+0{,}00093]$,
$p<0{,}001$), και μάλιστα σε \emph{κάθε} μία από τις 6 λίγκες, συμπεριλαμβανομένης της
Σκωτίας, όπου η διαφορά είναι η μικρότερη ($0{,}93590$ έναντι $0{,}93635$).
Πιθανή εξήγηση: κάθε λίγκα έχει μόνο 2.231--3.800 αγώνες (πολύ λίγο ώστε ο online
αλγόριθμος να συγκλίνει καλά μόνος του), και οι bookmakers που εμπιστεύεται περισσότερο ο
κάθε ειδικός διαφέρουν ελάχιστα ανά λίγκα: ο PSC είναι πρώτος σε 4 από τις 6 λίγκες και ο
B365 στις άλλες 2, δηλαδή οι ειδικοί καταλήγουν να «ανακαλύπτουν» σχεδόν τους ίδιους
έμπιστους experts με το ενιαίο μοντέλο, αντί να βρίσκουν κάτι τοπικό.

\subsection{Value betting}
\label{sec:core-algos-value}

Οικονομικός έλεγχος υγείας: θα έβγαζε πραγματικά λεφτά ένας bettor, μετά την αφαίρεση του
περιθωρίου (vig) του bookmaker; Το \file{value\_betting.py} χρησιμοποιεί σχεδιασμό
leave-one-bookmaker-out (ξαναπροσαρμόζει το OGD \emph{εξαιρώντας} τον bookmaker έναντι του
οποίου στοιχηματίζει) και προσομοίωση Kelly (κριτήριο Kelly -- βλ.\ ενότητα «Κριτήριο Kelly
(value betting)»).

\paragraph{Το panel εκπαίδευσης ταιριάζει στη φάση του στόχου.} Δεν αρκεί να φύγει η στήλη
του στόχου. Για \textbf{opening} στόχο πρέπει να φύγει ολόκληρη η closing πλευρά του panel:
καμία closing τιμή -- ούτε καν του ίδιου bookmaker -- δεν υπάρχει ακόμα τη στιγμή που είναι
διαθέσιμη μόνο η opening τιμή ενός αγώνα, αφού οι closing τιμές διαμορφώνονται όλες περίπου
ταυτόχρονα λίγο πριν το kickoff, ανεξάρτητα πότε άνοιξε η καθεμία. Ένα panel που τις αφήνει
μέσα στοιχηματίζει γνωρίζοντας πού θα κλείσει η αγορά -- look-ahead bias που παράγει τελικές
τράπεζες τάξης $10^6$--$10^9$, καθαρό τεχνούργημα κατασκευής. Με το panel περιορισμένο σωστά
σε opening-only, καμία τέτοια άκρη δεν επιβιώνει (π.χ.\ έναντι BW: $-0{,}00010$, $p=0{,}169$
raw· $+0{,}00004$, $p=0{,}439$ βαθμονομημένο).

Για \textbf{closing} στόχο δεν υπάρχει αντίστοιχος αιτιώδης περιορισμός, αλλά το
closing-only panel είναι ούτως ή άλλως το καλύτερο (\S\ref{sec:bestconfig}), και το
χρησιμοποιούμε. Η \texttt{run\_training\_panel\_sweep} τρέχει και τους τρεις τύπους panel
(closing-trained / mixed / opening-trained) $\times$ 2 κανονικοποιήσεις $\times$ κάθε στόχο,
108 γραμμές συνολικά, ώστε κάθε συνδυασμός να είναι ορατός.

\paragraph{Πλήρης σάρωση, 11 bookmakers.} Κάθε bookmaker με κάλυψη $\geq$50\% και πολλαπλές
σεζόν, βαθμονομημένες πιθανότητες, Kelly $\tfrac14$:

\begin{table}[htbp]
\centering
\small
\begin{tabular}{@{}llrrrl@{}}
\toprule
\textbf{Στόχος} & \textbf{Φάση} & \textbf{Στοιχήματα} & \textbf{Τράπεζα} &
\textbf{log-growth} & \textbf{$p$} \\
\midrule
WHC     & closing & 7.253  & $84{,}24\times$ & $+0{,}00061$ & $<0{,}001$ \\
IW      & opening & 7.443  & $7{,}89\times$  & $+0{,}00028$ & $0{,}008$ \\
BWC     & closing & 6.920  & $3{,}09\times$  & $+0{,}00016$ & $0{,}042$ \\
\midrule
VC\_BVC & closing & 4.561  & $1{,}58\times$  & $+0{,}00010$ & $0{,}339$ \\
WH      & opening & 5.931  & $1{,}38\times$  & $+0{,}00005$ & $0{,}545$ \\
BW      & opening & 7.841  & $1{,}09\times$  & $+0{,}00001$ & $0{,}820$ \\
B365C   & closing & 7.643  & $1{,}05\times$  & $+0{,}00001$ & $0{,}968$ \\
VC\_BV  & opening & 8.083  & $0{,}60\times$  & $-0{,}00006$ & $0{,}245$ \\
PSC     & closing & 19.497 & $0{,}54\times$  & $-0{,}00003$ & $0{,}368$ \\
B365    & opening & 10.829 & $0{,}41\times$  & $-0{,}00008$ & $0{,}101$ \\
PS      & opening & 26.675 & $0{,}26\times$  & $-0{,}00005$ & $0{,}055$ \\
\bottomrule
\end{tabular}
\caption{Value betting με panel ταιριασμένο στη φάση του στόχου. Η στήλη σύγκρισης είναι το
μέσο log-growth ανά στοίχημα· οι τελικές τράπεζες δεν συγκρίνονται μεταξύ γραμμών, αφού ο
αριθμός στοιχημάτων διαφέρει κατά τάξη μεγέθους. Πηγή:
\file{results/value\_betting\_training\_panel\_sweep.csv}.}
\label{tab:valuebetting}
\end{table}

\textbf{Τρεις από τους έντεκα} δίνουν σημαντική θετική άκρη (WHC, IW, BWC) και
\textbf{κανένας δεν δίνει σημαντική ζημιά}. Ο PSC -- ο \#1 προγνώστης συνολικά
(\S\ref{sec:core-algos}) -- βγαίνει ισοπαλία ($p=0{,}368$), το ίδιο αποτέλεσμα που δίνει και
η σύγκριση πρόβλεψης στη διαμόρφωση αναφοράς (\S\ref{sec:bestconfig}), από ανεξάρτητο δρόμο:
δύο διαφορετικές μετρικές λένε ότι το closing-only μείγμα φτάνει την Pinnacle χωρίς να την
ξεπερνά.

\paragraph{Η αντιστοίχιση φάσης δεν είναι λεπτομέρεια.} Η σάρωση των 108 γραμμών επιτρέπει να
μετρηθεί άμεσα:
\begin{itemize}
  \item Με \textbf{mixed} panel αντί για closing-only, το αποτέλεσμα χειροτερεύει σε
        \emph{κάθε} έναν από τους πέντε closing στόχους, και ο PSC γίνεται σημαντική ζημιά
        ($-0{,}00021$, $p<0{,}001$) αντί για ισοπαλία.
  \item Με \textbf{opening-trained} μείγμα εναντίον closing στόχων, η ζημιά είναι σημαντική
        σε \emph{όλους} ($-0{,}00044$ έως $-0{,}00071$, όλα $p<0{,}001$). Παλιότερη
        πληροφορία εναντίον νεότερης τιμής χάνει πάντα -- το οικονομικό ανάλογο της κλίσης
        του \S\ref{sec:infoarrival}. Αυτή η κατεύθυνση δεν έχει πρόβλημα αιτιότητας, απλώς
        χάνει.
\end{itemize}

\begin{figure}[htbp]
  \centering
  \includegraphics[width=0.92\linewidth]{value_betting_all_bookmakers.png}
  \caption{Σάρωση άκρης στοιχηματισμού: μέση λογαριθμική ανάπτυξη ανά στοίχημα, raw
  (κενός κύκλος) έναντι calibrated (γεμάτος), leave-one-out ανά bookmaker.
  Αστερίσκος = σημαντικό στο 95\%.}
\end{figure}

Η σωστή διατύπωση δεν είναι «καμία άκρη έναντι σκληρού, μεγάλη έναντι αδύναμου» -- είναι:
\textbf{η ύπαρξη και το μέγεθος της άκρης εξαρτώνται από τον συγκεκριμένο bookmaker}. Σε κάθε
περίπτωση, όρια λογαριασμού/ρευστότητας σε πραγματικούς λογαριασμούς κερδισμένων παικτών δεν
μοντελοποιούνται -- η προσομοίωση δείχνει στατιστική ύπαρξη άκρης, όχι εγγυημένη, βιώσιμη
κερδοφορία στην πράξη.

Συμπληρωματικά, μια καθαρά περιγραφική δοκιμή στο \file{market\_devig\_comparisons.py}
(\texttt{run\_quality\_transfer\_correlation}) ελέγχει αν η «εξυπνάδα» ενός bookmaker στο
closing μεταφέρεται στο δικό του opening (13 ζεύγη, Pearson $r=-0{,}006$ σε όλα τα ζεύγη,
$r=-0{,}390$ σε 7 multi-season ζεύγη μόνο, Spearman $\rho \in [+0{,}18,\,+0{,}29]$) --
ασταθές πρόσημο, μικρό δείγμα, καμία αξιόπιστη σχέση.

\subsection{Ακρίβεια ανά κατηγορία αποτελέσματος}

Κάθε άλλη μετρική στο project μέσο-όρίζει πάνω και στα τρία ενδεχόμενα μαζί. Το
\file{outcome\_class\_accuracy.py} σπάει την ακρίβεια ξεχωριστά ανά Νίκη
Γηπεδούχου/Ισοπαλία/Νίκη Φιλοξενούμενου. Εύρημα, καλά τεκμηριωμένο στη βιβλιογραφία
πρόβλεψης ποδοσφαίρου και επιβεβαιωμένο εδώ: οι \textbf{ισοπαλίες είναι ουσιαστικά
απρόβλεπτες} ως argmax επιλογή για κάθε αλγόριθμο/bookmaker: το καλύτερο (PSC) πετυχαίνει
το $0{,}695\%$ των πραγματικών ισοπαλιών και το OGD το $0{,}299\%$, παρότι οι ισοπαλίες
αποτελούν $\sim$26\% των αγώνων. Για σύγκριση, οι νίκες γηπεδούχου προβλέπονται σωστά με
ακρίβεια έως $82{,}8\%$ (BW) και οι νίκες φιλοξενούμενου έως $49{,}3\%$ (PSC).

\subsection{Χρονικές τάσεις}
\label{sec:temporal}

Ωριμάζει η αγορά με τον χρόνο; Το \file{temporal\_trends.py} ελέγχει log-loss, ECE και
overround ανά σεζόν, με το επίσημο τεστ τάσης Mann--Kendall
(\S\ref{sec:bootstrap-theory}) αντί για απλή οπτική επιθεώρηση μιας πτωτικής γραμμής.
Εύρημα: \textbf{καμία σημαντική τάση} στην ακρίβεια της αγοράς (log-loss/ECE) στις 10
σεζόν (όλα τα $p>0{,}37$), αλλά το \textbf{overround αυξάνεται σημαντικά} (μέσος όρος
bookmakers: $4{,}42\%$ το 2016/17 $\to 6{,}0\%$ το 2025/26, $p=0{,}0007$) --
το αντίθετο μιας αφήγησης «ωριμάζουσας, πιο ανταγωνιστικής αγοράς». Αυτό είναι επίσης
πιθανή εξήγηση για το γιατί ένα ενιαίο global fit βαθμονόμησης δεν προλαβαίνει να
προσαρμοστεί καλά σε όλες τις σεζόν (\S\ref{sec:calib-theory}).

\subsection{Σύγκριση παραθύρου pooling (5ετία έναντι 10ετίας)}
\label{sec:pooling}

Το project ξεκίνησε σε 5 σεζόν/10 λίγκες και επεκτάθηκε σε 10 σεζόν/22 λίγκες· αρκετά
ευρήματα άλλαξαν ανάμεσα στις δύο κλίμακες. Το \file{pooling\_window\_comparison.py}
ξανατρέχει το ίδιο pipeline κατάταξης (\file{final\_ranking.py}) και μια παρόμοια μπαταρία
σημαντικότητας \textbf{από την αρχή}, ανεξάρτητα, σε δύο χρονικά παράθυρα (2021/22--2025/26
έναντι 2016/17--2025/26 πάνω στο ίδιο διευρυμένο 22-λίγκο panel), απομονώνοντας καθαρά την
επίδραση του βάθους δεδομένων από την επέκταση των λιγκών.

\subsection{Opening έναντι Closing αποδόσεων}
\label{sec:openclose-results}

Χτίστηκε σε αυτή τη συνεδρία (\file{market\_devig\_comparisons.py},
\texttt{run\_opening\_vs\_closing}), απαντώντας: κάνει η τελική,
πιο ενημερωμένη τιμή της αγοράς (closing) καλύτερο μείγμα από την πρώτη (opening); Το panel
των 26 experts χωρίζεται καθαρά σε 13/13 (\S\ref{sec:openclose}). Ξανατρέχονται οι 4
αλγόριθμοι πυρήνα ανεξάρτητα σε \textbf{τρία} panels: μόνο opening, μόνο closing, και το
πλήρες μικτό panel των 26 -- \textbf{χωρίς} νέο αλγοριθμικό κώδικα, απλά περιορίζοντας το
panel (οι συναρτήσεις sleeping-experts ήδη δέχονται αυθαίρετο υποσύνολο experts).

Ευρήματα, στατιστικά επιβεβαιωμένα (moving block bootstrap, $p\approx 0$ παντού) και για
τους 4 αλγορίθμους:
\begin{itemize}
  \item \textbf{Closing $<$ πλήρες panel $<$ opening} (χαμηλότερο log-loss = καλύτερο) --
        η διάταξη ισχύει με στατιστική βεβαιότητα, όχι μόνο ως τυχαία διαφορά σημειακών
        εκτιμήσεων.
  \item Το closing-only μείγμα δεν είναι απλώς καλύτερο από το opening-only -- είναι
        καλύτερο κι από το \textbf{πλήρες μικτό panel} των 26 (π.χ.\ OGD: closing-only
        0{,}99722 έναντι πλήρους 0{,}99745).
  \item Το opening-only \textbf{υποαποδίδει σημαντικά} έναντι του πλήρους panel επίσης --
        δηλαδή οι opening bookmakers λειτουργούν ως καθαρό «βάρος» στο μείγμα, όχι ως
        χρήσιμη διαφοροποίηση.
  \item Σε ενιαία κατάταξη 38 εγγραφών (12 παραλλαγές αλγορίθμων $\times$ 26 μεμονωμένοι
        bookmakers, Σχήμα~\ref{fig:ranking}), το \textbf{OGD (closing) είναι ο καλύτερος
        αλγόριθμος συνολικά, στη θέση \#3} -- πίσω μόνο από το PSC και μια σειρά
        Betfair-closing μονής σεζόν -- μπροστά από κάθε άλλο bookmaker και κάθε
        opening/full παραλλαγή αλγορίθμου.
\end{itemize}

\begin{figure}[htbp]
  \centering
  \includegraphics[width=0.95\linewidth]{opening_vs_closing_ranking.png}
  \caption{Πλήρης κατάταξη: 12 παραλλαγές αλγορίθμων (opening/closing/full panel) + 26
  μεμονωμένοι bookmakers, ταξινομημένοι κατά log-loss.}
  \label{fig:ranking}
\end{figure}

\subsubsection{Πού ζει το προβάδισμα των closing αποδόσεων}
\label{sec:infoarrival}

Το «τα closing είναι καλύτερα» είναι μέσος όρος, και ένας μέσος όρος μπορεί να κρύβει ότι το
φαινόμενο συμβαίνει μόνο σε ένα υποσύνολο. Το \file{information\_arrival.py} στρωματοποιεί
τους αγώνες κατά δύο δείκτες άφιξης πληροφορίας και ξαναμετρά τη διαφορά closing μείον
opening μέσα σε κάθε στρώμα:

\begin{itemize}
  \item \textbf{Κίνηση γραμμής}: η συνολική μεταβολή $\sum_k |p_k^{\text{close}} -
        p_k^{\text{open}}|$ ανά αγώνα, δηλαδή πόσο μετακινήθηκε η αγορά.
  \item \textbf{Ανωμαλία $z$ του Shin}: πόσο αποκλίνει η εκτιμώμενη αναλογία «insider»
        στοιχηματιστών του αγώνα από τη συνήθη τιμή της (\S\ref{sec:shin}).
\end{itemize}

Στην κίνηση γραμμής η κλίση είναι δραματική και μονοτονική: στο χαμηλότερο πεμπτημόριο η
διαφορά είναι $-0{,}00044$ και \emph{δεν} είναι σημαντική ($p=0{,}052$), στο υψηλότερο
$-0{,}01158$ ($p<0{,}001$) --- λόγος $26\times$ από άκρη σε άκρη. Στην ανωμαλία $z$ η κλίση
υπάρχει αλλά είναι πολύ πιο ήπια ($-0{,}00295$ έως $-0{,}00498$, σημαντική σε όλα τα
στρώματα).

Η ερμηνεία είναι σημαντική για το πώς διαβάζεται όλο το \S\ref{sec:openclose-results}: η
απόδοση κλεισίματος δεν είναι καλύτερη επειδή είναι μεταγενέστερη, αλλά επειδή --- σε
ορισμένους αγώνες --- έχει απορροφήσει πληροφορία που έφτασε στο ενδιάμεσο. Όπου η αγορά δεν
κινήθηκε, η τιμή κλεισίματος δεν ξέρει μετρήσιμα περισσότερα από την αρχική. Αυτό συνδέεται
άμεσα με τη βιβλιογραφία μικροδομής της αγοράς, όπου η κίνηση της τιμής είναι ακριβώς ο
μηχανισμός με τον οποίο η ιδιωτική πληροφορία ενσωματώνεται στην τιμή (Kyle 1985; Glosten
\& Milgrom 1985).

\subsection{Εναλλακτική μέθοδος de-vig (Shin)}
\label{sec:shin}

Όλο το project de-vigάρει τις ωμές αποδόσεις με proportional normalization
($p_i = \pi_i / \sum_j \pi_j$, όπου $\pi_i = 1/\text{odds}_i$) -- μια επιλογή που ποτέ δεν
είχε ελεγχθεί έναντι εναλλακτικής. Το \file{data\_processer.py} χτίζει, στο ίδιο run,
και ένα δεύτερο dataset με τη μέθοδο του \textbf{Shin (1992)}: μοντελοποιεί μια αναλογία
$z$ «insider» στοιχηματιστών που ποντάρουν με βεβαιότητα στο πραγματικό αποτέλεσμα, και
λύνει (μέσω bisection, έναν αγώνα τη φορά) για το $z$ και τις δίκαιες πιθανότητες
ταυτόχρονα -- αντί να μοιράζει το περιθώριο του bookmaker ομοιόμορφα, το μεταφέρει από το
outsider προς το φαβορί, η γνωστή διόρθωση του favorite-longshot bias.

Τα ευρήματα (\file{market\_devig\_comparisons.py}'s \texttt{run\_shin\_vs\_basic\_comparison}
και \texttt{run\_shin\_vs\_basic\_calibration}, \file{value\_betting.py}'s
\texttt{run\_shin\_vs\_basic}):
\begin{itemize}
  \item Η \textbf{κατάταξη δεν αλλάζει}: ο PSC παραμένει \#1, ίδια σειρά αλγορίθμων/
        bookmakers, είτε με proportional είτε με Shin de-vig.
  \item Στο \textbf{raw} στάδιο, το Shin είναι σημαντικά καλύτερο (log-loss \emph{και}
        Brier) σε όλες τις 11 Tier-A οντότητες, $p<0{,}0001$ παντού.
  \item Η ποιότητα βαθμονόμησης (ECE) βελτιώνεται δραματικά -- μειώνεται $\sim$35--57\%
        σε κάθε σειρά (π.χ.\ OGD: $0{,}0091 \to 0{,}0043$).
  \item Μετά τη βαθμονόμηση όμως, το πλεονέκτημα συρρικνώνεται ή εξαφανίζεται· και,
        αντίστροφα, το ίδιο το calibration παύει να είναι σημαντικό για τα περισσότερα
        series \emph{μέσα} στο Shin (μόνο 2 από τα 11 παραμένουν σημαντικά, έναντι 11/11
        στο proportional). Ερμηνεία: το Shin και το temperature scaling διορθώνουν το
        \emph{ίδιο} favorite-longshot bias, με διαφορετικό μηχανισμό (per-match
        cross-sectional το ένα, online πάνω σε όλο το ιστορικό το άλλο) -- μόλις γίνει η
        μία διόρθωση, η άλλη γίνεται σε μεγάλο βαθμό περιττή.
  \item Το ίδιο μοτίβο επιβεβαιώνεται και στην οικονομική προσομοίωση (Kelly betting, για
        τους 2 χειρότερους αλγορίθμους FTRL/Uniform average): σημαντικά καλύτερο raw
        παντού, αλλά μόνο σε 2 από τις 8 συγκρίσεις μετά τη βαθμονόμηση.
\end{itemize}

\begin{figure}[htbp]
  \centering
  \includegraphics[width=0.92\linewidth]{shin_vs_basic_calibration_reliability.png}
  \caption{Calibration gap: proportional normalization (αριστερά) έναντι Shin (δεξιά) --
  το εύρος συρρικνώνεται δραματικά, οπτική επιβεβαίωση της διόρθωσης favorite-longshot
  bias.}
\end{figure}

\textbf{Συμπέρασμα:} το βασικό εύρημα του project (το μείγμα ξεπερνά κάθε μεμονωμένο
bookmaker εκτός του καλύτερου sharp bookmaker· η βαθμονόμηση βοηθάει) \textbf{δεν είναι
artifact} της επιλογής de-vig -- επιβεβαιώνεται ανεξάρτητα με δύο μεθόδους.

\subsection{Η καλύτερη διαμόρφωση: στοίβαξη των μετρημένων βελτιώσεων}
\label{sec:bestconfig}

Τέσσερα πράγματα είχαν μετρηθεί χωριστά και ποτέ μαζί: το διορθωμένο βήμα του OGD
(\S\ref{sec:steptuning}), το closing-only panel που ξεπερνά το πλήρες
(\S\ref{sec:openclose-results}), το Shin de-vig που ξεπερνά την proportional normalization
στο raw στάδιο (\S\ref{sec:shin}), και η βαθμονόμηση (\S\ref{sec:calib-theory}). Το
pipeline έτρεχε με το πλήρες panel και proportional normalization, δηλαδή άφηνε δύο από τα
τέσσερα έξω. Το \file{best\_configuration.py} τρέχει τη στοίβαξη από άκρη σε άκρη και ρωτά το
ένα ερώτημα που μπορεί να αλλάξει την κεντρική αφήγηση: \emph{υπάρχει διαμόρφωση που φτάνει
τον PSC σε μια έντιμη, ισότιμη σύγκριση;}

Δύο προφυλάξεις, και οι δύο επειδή το project έχει πληρώσει την απουσία τους στο παρελθόν:

\begin{enumerate}
  \item \textbf{Ίδιοι αγώνες για όλους.} Κάθε διαμόρφωση βαθμολογείται στην τομή της κάλυψης
        όλων των διαμορφώσεων με του PSC --- 71.818 γύροι. Τα panel έχουν διαφορετική κάλυψη,
        και ένα log-loss υπολογισμένο πάνω σε διαφορετικό σύνολο αγώνων δεν είναι συγκρίσιμο
        μέγεθος. Η \emph{εκπαίδευση} χρησιμοποιεί κανονικά όλο το διαθέσιμο εύρος κάθε
        διαμόρφωσης· μόνο η βαθμολόγηση περιορίζεται.
  \item \textbf{Held-out ουρά.} Το να διαλέγεις διαμόρφωση κοιτώντας αυτά τα νούμερα
        \emph{είναι} επιλογή, με το αντίστοιχο κόστος data snooping. Όλα αναφέρονται δύο
        φορές: πλήρες δείγμα (περιγραφικά) και χρονική ουρά 25\% που η επιλογή δεν είδε ποτέ.
\end{enumerate}

\begin{table}[htbp]
\centering
\small
\begin{tabular}{@{}lrrr@{}}
\toprule
\textbf{Διαμόρφωση} & \textbf{πλήρες δείγμα} & \textbf{train 75\%} & \textbf{held-out 25\%} \\
\midrule
PSC μόνος του                            & 0,99685 & 0,99746 & 0,99502 \\
\textbf{A: closing-only + Shin}          & \textbf{0,99692} & 0,99754 & 0,99507 \\
closing-only, proportional               & 0,99693 & 0,99756 & 0,99506 \\
baseline: πλήρες panel, proportional     & 0,99728 & 0,99801 & 0,99508 \\
πλήρες panel, Shin                       & 0,99731 & 0,99804 & 0,99512 \\
B: A χωρίς bookmakers μονής σεζόν        & 0,99692 & 0,99754 & 0,99509 \\
\bottomrule
\end{tabular}
\caption{Βαθμονομημένο log-loss του OGD ανά διαμόρφωση, ίδιοι 71.818 αγώνες σε κάθε γραμμή.
Πηγή: \file{results/best\_configuration.csv}.}
\label{tab:bestconfig}
\end{table}

Τα ζευγαρωτά τεστ (moving block bootstrap πάνω στους κοινούς γύρους) δείχνουν καθαρά ποιος
παράγοντας δουλεύει και ποιος όχι:

\begin{itemize}
  \item \textbf{Το panel κάνει όλη τη δουλειά.} Η διαμόρφωση A ξεπερνά το baseline κατά
        $-0{,}00036$ ($p<0{,}001$)· το closing-only με proportional normalization πετυχαίνει
        σχεδόν το ίδιο ($-0{,}00034$, $p<0{,}001$).
  \item \textbf{Το Shin μόνο του δεν προσφέρει τίποτα μετρήσιμο} πάνω στο πλήρες panel
        ($+0{,}00004$, $p=0{,}181$) --- απόλυτα συνεπές με το εύρημα του \S\ref{sec:shin}
        ότι το Shin και το temperature scaling διορθώνουν το ίδιο bias, οπότε το δεύτερο
        απορροφά το πρώτο. Η A είναι λοιπόν η καλύτερη διαμόρφωση, αλλά το «$+$Shin» σε αυτό
        το όνομα είναι διακοσμητικό.
  \item \textbf{Έναντι του PSC}: το baseline χάνει σημαντικά ($+0{,}00043$, $p<0{,}001$),
        ενώ η A \emph{ισοπαλεί} ($+0{,}00007$, $p=0{,}234$) --- όπως και το closing-only με
        proportional ($+0{,}00009$, $p=0{,}143$).
  \item Η διαμόρφωση B (η A χωρίς τους bookmakers μονής σεζόν) είναι οριακά χειρότερη από την
        A στην held-out ουρά και απορρίπτεται: το φιλτράρισμα ποιότητας δεν προσθέτει τίποτα
        πάνω από το ίδιο το panel.
\end{itemize}

\textbf{Τι σημαίνει και τι δεν σημαίνει αυτό.} Σημαίνει ότι η διατύπωση «ο PSC μας ξεπερνά
σημαντικά» ισχύει για το πλήρες panel ($p<0{,}001$) αλλά όχι για τη διαμόρφωση αναφοράς
($p=0{,}234$). \emph{Δεν} σημαίνει ότι το μείγμα ξεπερνά τον PSC --- το σημειακό
πρόσημο παραμένει υπέρ του PSC, και στην held-out ουρά και οι έξι γραμμές του
Πίνακα~\ref{tab:bestconfig} μαζεύονται σε λιγότερο από $0{,}0001$ μεταξύ τους. Η ουρά
επιβεβαιώνει την ισοπαλία, δεν αναδεικνύει νικητή, και θα ήταν ανέντιμο να παρουσιαστεί ως
νίκη. Τέλος, σε κάθε closing διαμόρφωση ο PSC βρίσκεται \emph{μέσα} στο μείγμα: το «φτάνουμε
τον PSC» διαβάζεται «το μείγμα φτάνει το καλύτερο μέλος του», όχι «τον φτάνουμε από έξω» ---
αυτό το τελευταίο το ελέγχει το \S\ref{sec:nopinnacle}, με το αποτέλεσμα που περιγράφεται
εκεί.

% =====================================================================
\section{Ζητήματα ποιότητας δεδομένων που εντοπίστηκαν}
% =====================================================================

Μέσω ελέγχου έναντι της ίδιας τεκμηρίωσης του football-data.co.uk:
\begin{itemize}
  \item Rebrand VC$\to$BetVictor: συγχώνευση σε μία ταυτότητα (\S\ref{sec:dataproc}).
  \item Αποκλεισμός στηλών συγκεντρωτικών αποδόσεων Betbrain/Max/Avg (ορατό μόνο όταν το
        dataset επεκτάθηκε πίσω από το $\sim$2021).
  \item Bookmakers ενεργοί σε \textbf{μία μόνο σεζόν} επισημαίνονται ρητά (πεδίο
        \var{single\_season} στο \file{results\_table.csv}) -- η φαινομενική τους απόδοση
        είναι συγχυσμένη με τη δυσκολία εκείνης συγκεκριμένα της σεζόν, όχι αξιόπιστη
        σύγκριση δεξιοτήτων.
\end{itemize}

\subsection{Πώς διαβάζεται το βάρος ενός κοιμισμένου expert}
\label{sec:bugfix}

Οι τρεις αλγόριθμοι γράφουν στον πίνακα βαρών $W[t,i]$ \textbf{μόνο} για τους experts που
είναι ξύπνιοι στον γύρο $t$ (οι κοιμισμένες συντεταγμένες της γραμμής μένουν μηδέν, εξ
ορισμού της αναγωγής Sleeping Experts, \S\ref{sec:sleeping}). Συνεπώς η \emph{τελευταία
γραμμή} $W[-1]$ \textbf{δεν} είναι το τελικό βάρος κάθε expert: είναι ακριβώς μηδέν για
όποιον δεν έδωσε απόδοση στον τελευταίο αγώνα του dataset -- 10 από τους 26 bookmakers εδώ,
μεταξύ τους ο PSC, ο καλύτερος μεμονωμένος προγνώστης του συνόλου. Διαβασμένο ως «τελικό
βάρος», αυτό το μηδέν σημαίνει «δεν κάλυψε τον τελευταίο αγώνα», όχι «ο αλγόριθμος έμαθε να
τον αγνοεί».

Η ποσότητα που αναφέρεται παντού στην εργασία (\var{last\_awake\_weight} στο
\file{sleeping\_experts.py}) είναι το βάρος που κρατούσε ο κάθε expert \textbf{την τελευταία
φορά που ήταν ξύπνιος}. Κάθε τιμή είναι μερίδιο του συνόλου των ξύπνιων \emph{εκείνου} του
γύρου, άρα ερμηνεύσιμη μία-μία («πόσο της πρόβλεψης εκείνου του γύρου ήταν δική του, την
τελευταία φορά που συμμετείχε»), αλλά οι τιμές δεν αθροίζουν σε $1$ μεταξύ experts. Η
διάκριση αφορά μόνο την ανάγνωση της κατάστασης, όχι τους ίδιους τους αλγορίθμους: καμία
μετρική ποιότητας πρόβλεψης (log-loss, Brier, RPS, accuracy) δεν εξαρτάται από αυτήν. Είναι
όμως ουσιώδης για τους πίνακες βαρών του \file{fixed\_share.py} και του
\file{contextual\_experts.py}, όπου αφορά ακριβώς τους bookmakers μονής σεζόν και τον PSC --
δηλαδή τις σειρές που τα δύο αυτά αρχεία υπάρχουν για να εξετάσουν.

% =====================================================================
\section{Γνωστοί περιορισμοί}
% =====================================================================

\begin{itemize}
  \item Η συγχώνευση VC/BV συνάγεται από τη μη-επικαλυπτόμενη παρουσία ανά σεζόν, όχι
        επιβεβαιωμένη από επίσημη πηγή -- πιθανόν να υπάρχουν κι άλλα μη εντοπισμένα
        rebrands, ιδίως τώρα με 22 λίγκες.
  \item Το log-loss είναι η κύρια μετρική στο μεγαλύτερο μέρος του project· το Brier
        αναλύεται συστηματικά μόνο στη σύγκριση de-vig μεθόδων (ενότητα «Εναλλακτική
        μέθοδος de-vig (Shin)»)· το RPS υπολογίζεται παντού αλλά δεν αναλύεται
        συστηματικά πουθενά.
  \item Η προσομοίωση value betting αγνοεί όρια λογαριασμού/ρευστότητας.
  \item Ένα leave-one-bookmaker-out σχέδιο εισάγει look-ahead bias αν το panel εξαίρεσης
        δεν περιοριστεί στην πληροφορία που θα ήταν πραγματικά διαθέσιμη τη στιγμή του
        στοιχήματος. Ο περιορισμός εφαρμόζεται ρητά ανά φάση αγοράς (ενότητα «Value
        betting»), αλλά δεν αποκλείεται να υπάρχουν άλλα, λιγότερο προφανή κανάλια
        διαρροής στο project.
  \item Οι αιτιώδεις εξηγήσεις για κάποια ευρήματα (π.χ.\ γιατί αυξάνεται το overround)
        είναι εύλογες, όχι αποδεδειγμένες.
\end{itemize}

% =====================================================================
\section{Ανοιχτά επόμενα βήματα}
% =====================================================================

\begin{itemize}
  \item Δεύτερο άθλημα (π.χ.\ μπάσκετ) για έλεγχο γενίκευσης του πλαισίου OCO πέρα από το
        ποδόσφαιρο.
  \item Πραγματικό baseline Bayesian Model Averaging για εμπειρική σύγκριση (συζητείται
        θεωρητικά, ποτέ δεν υλοποιήθηκε).
  \item Συστηματική αναζήτηση για περισσότερες περιπτώσεις rebrand/merge bookmakers σε
        όλες τις 22 λίγκες.
  \item Διερεύνηση του \emph{γιατί} αυξάνεται το overround (θα χρειαστεί δεδομένα εκτός
        football-data.co.uk).
\end{itemize}

% =====================================================================
\appendix
\section{Πλήρης χάρτης αρχείων}
\label{app:filemap}
% =====================================================================

Πίνακας κάθε αρχείου στο \file{scripts/}: ρόλος, βασικές εξαρτήσεις, και τι παράγει στο
\file{results/} ή \file{docs/}. Σειρά: η σειρά εκτέλεσης/εξάρτησης (βλ.\ και το διάγραμμα
της \S\ref{sec:architecture}).

{\small
\begin{longtable}{@{}>{\raggedright\arraybackslash}p{3.7cm} >{\raggedright\arraybackslash}p{4.7cm} >{\raggedright\arraybackslash}p{6.0cm}@{}}
\toprule
\textbf{Αρχείο} & \textbf{Ρόλος} & \textbf{Έξοδοι} \\
\midrule
\endfirsthead
\toprule
\textbf{Αρχείο} & \textbf{Ρόλος} & \textbf{Έξοδοι} \\
\midrule
\endhead
\file{download\_data.ps1} &
  Κατεβάζει τα ωμά CSV από το football-data.co.uk (22 λίγκες $\times$ 10 σεζόν). &
  \file{data/raw/*.csv}, \file{data/download\_log.csv} \\[4pt]
\file{data\_processer.py} &
  Εντοπισμός τριάδων αποδόσεων (hardcoded λίστα bookmakers), de-vig -- ΚΑΙ proportional
  ΚΑΙ Shin normalization στο ίδιο run, merges/εξαιρέσεις. Αντικαθιστά το παλιότερο δίπτυχο
  \file{process\_data.py} + \file{process\_data\_shin.py}
  (\file{scripts/deprecated/}, byte-/τιμή-ταυτόσημη έξοδος επαληθευμένη). &
  \file{data/processed/odds\_long.csv}, \file{odds\_long\_shin.csv} \\[4pt]
\file{sleeping\_experts.py} &
  \textbf{Θεμέλιο -- μόνο algorithms/data-loading, χωρίς \var{main()}.} Γενικεύει
  Hedge/OGD/FTRL σε ολόκληρο το panel (26 experts, μερική κάλυψη) μέσω αναγωγής Sleeping
  Experts. \var{load\_full\_universe()}, \var{run\_*\_sleeping()}, \var{evaluate()}. &
  (καμία -- library αρχείο) \\[4pt]
\file{sleeping\_experts\_experiment.py} &
  Τρέχει τους αλγορίθμους του \file{sleeping\_experts.py} πάνω στο πλήρες panel. &
  \file{results/results\_table.csv}, \file{sleeping\_experts\_comparison.png} \\[4pt]
\file{calibration\_analysis.py} &
  Διαγνωστικό: reliability diagram, ECE. Εντοπίζει το favorite-longshot bias, δεν το
  διορθώνει. &
  \file{results/calibration\_table.csv}, \file{calibration\_reliability.png} \\[4pt]
\file{calibration\_correction.py} &
  Online temperature scaling (\var{temperature\_calibrate}, \var{select\_lr\_on\_train}),
  train/val πρωτόκολλο. &
  \file{results/calibration\_correction\_table.csv}, \file{calibration\_correction.png} \\[4pt]
\file{final\_ranking.py} &
  Ενιαία επίσημη κατάταξη Tier-A (κάλυψη $\geq$70\%): raw ΚΑΙ calibrated, ίδιο πρωτόκολλο,
  άμεσα συγκρίσιμα. \var{calibrate\_full\_history()} εισάγεται από 6 άλλα αρχεία. &
  \file{results/final\_ranking\_table.csv}, \file{final\_ranking.png} \\[4pt]
\file{significance\_test.py} &
  Moving block bootstrap, μπαταρία 12 συγκρίσεων. \var{moving\_block\_bootstrap\_test},
  \var{paired\_diff} επαναχρησιμοποιούνται από \file{value\_betting.py},
  \file{value\_betting\_online\_kelly.py}, \file{contextual\_experts.py},
  \file{pooling\_window\_comparison.py}, \file{market\_devig\_comparisons.py}. &
  \file{results/significance\_test\_table.csv} \\[4pt]
\file{fixed\_share.py} &
  Διόρθωση «παγωμένου βάρους» σπάνιων bookmakers μέσω διαρροής προς το ομοιόμορφο. &
  \file{results/fixed\_share\_table.csv}, \file{fixed\_share\_weights.csv} \\[4pt]
\file{contextual\_experts.py} &
  Βοηθάει η εξειδίκευση βαρών ανά λίγκα; (Όχι, εκτός Σκωτίας.) &
  \file{results/contextual\_experts\_table.csv}, \file{contextual\_experts.png} \\[4pt]
\file{value\_betting.py} &
  \emph{(Ενοποιημένο αρχείο, 6 πειράματα.)} Προσομοίωση Kelly leave-one-bookmaker-out: θα
  έβγαζε λεφτά ένας πραγματικός bettor; Ο σημειακός έλεγχος (B365C/BW), η πλήρης
  σάρωση 11 bookmakers, το closing-on-closing και για τα δύο normalizations, ο έλεγχος
  αιτιότητας του panel έναντι BW, το πλήρες grid 108 γραμμών, και το Shin-έναντι-basic value
  betting για τους 2 χειρότερους αλγορίθμους -- όλα σε ένα αρχείο, με μία γενικευμένη
  συνάρτηση προσαρμογής (\var{fit\_leave\_out}) αντί για 5 σχεδόν πανομοιότυπα αντίγραφα. &
  \file{results/value\_betting\_table.csv}, \file{value\_betting.png},
  \file{value\_betting\_all\_bookmakers\_table.csv}, \file{\_.png},
  \file{value\_betting\_closing\_trained\_table.csv}, \file{\_.png},
  \file{opening\_value\_betting\_corrected.csv}, \file{\_.png},
  \file{value\_betting\_training\_panel\_sweep.csv},
  \file{shin\_vs\_basic\_value\_betting\_table.csv}, \file{\_significance.csv}, \file{\_.png} \\[4pt]
\file{outcome\_class\_accuracy.py} &
  Ακρίβεια ξεχωριστά ανά Η/Χ/2 αντί για μέσο όρο. &
  \file{results/outcome\_class\_accuracy\_table.csv},
  \file{outcome\_class\_accuracy.png} \\[4pt]
\file{temporal\_trends.py} &
  Log-loss/overround/ECE ανά σεζόν + επίσημο τεστ τάσης Mann--Kendall. &
  \file{results/temporal\_trends\_table.csv}, \file{temporal\_trends\_mk\_test.csv},
  \file{temporal\_trends.png} \\[4pt]
\file{pooling\_window\_comparison.py} &
  Ξανατρέχει την κατάταξη/σημαντικότητα σε 2 χρονικά παράθυρα (5ετία/10ετία) από την αρχή
  το καθένα, για να απομονώσει την επίδραση βάθους δεδομένων. &
  \file{results/pooling\_window\_ranking\_table.csv},
  \file{pooling\_window\_significance\_table.csv}, \file{pooling\_window\_comparison.png} \\[4pt]
\file{calibration\_per\_season.py} &
  Επαναφορά βαθμονόμησης κάθε σεζόν (cold-start) αντί για ένα global fit. &
  \file{results/calibration\_per\_season\_table.csv}, \file{\_significance.csv},
  \file{calibration\_per\_season.png} \\[4pt]
\file{calibration\_warmstart.py} &
  \emph{(Χτίστηκε σε αυτή τη συνεδρία.)} Ενδιάμεσο variant: μηδενισμός μόνο του μετρητή
  γύρου, όχι του $s$, ανάμεσα σε σεζόν. &
  \file{results/calibration\_warmstart\_table.csv}, \file{\_significance.csv},
  \file{calibration\_warmstart.png} \\[4pt]
\file{market\_devig\_comparisons.py} &
  \emph{(Ενοποιημένο αρχείο, 4 πειράματα.)} (1) Χωρίζει το panel σε opening/closing μέσω
  της κατάληξης «C», ξανατρέχει τους 4 αλγορίθμους και στα δύο υποσύνολα + στο πλήρες
  panel, ενιαία κατάταξη με όλους τους bookmakers· (2) συσχέτιση closing-side με
  opening-side δεξιότητα ανά bookmaker· (3) ξανατρέχει αλγορίθμους + Tier-A bookmakers σε
  proportional έναντι Shin normalization (raw/calibrated log-loss+Brier, bootstrap
  σημαντικότητα)· (4) ECE / reliability diagram, proportional έναντι Shin. Μία κοινή
  συνάρτηση (\var{fit\_all\_algorithms}) αντί για 3 αντίγραφα του ίδιου υπολογισμού
  βαρών. &
  \file{results/opening\_vs\_closing\_table.csv}, \file{\_significance.csv},
  \file{opening\_vs\_closing.png}, \file{\_ranking.csv}, \file{\_ranking.png},
  \file{opening\_closing\_quality\_transfer.csv}, \file{\_.png},
  \file{shin\_vs\_basic\_table.csv}, \file{\_significance.csv},
  \file{\_calibration\_benefit.csv}, \file{shin\_vs\_basic\_comparison.png},
  \file{shin\_vs\_basic\_calibration\_table.csv}, \file{\_ece.csv},
  \file{shin\_vs\_basic\_calibration\_reliability.png} \\[4pt]
\file{value\_betting\_online\_kelly.py} &
  \emph{(Χτίστηκε σε αυτή τη συνεδρία.)} Το stake μαθαίνεται online (projected OGD πάνω
  στο $-\log(\text{wealth growth})$) αντί για τον κλειστού-τύπου τύπο Kelly -- τρίτο OCO
  στάδιο. Χάνει σχεδόν παντού από το κλασικό Kelly. &
  \file{results/value\_betting\_online\_kelly\_table.csv}, \file{\_significance.csv},
  \file{value\_betting\_online\_kelly.png} \\[4pt]
\bottomrule
\end{longtable}
}

\section{Πώς τρέχει το πλήρες pipeline}

Από μέσα στον φάκελο \file{scripts/} (οι σχετικές εισαγωγές μεταξύ αρχείων δεν επιλύονται
από τη ρίζα του project):
\begin{verbatim}
download_data.ps1
data_processer.py          # μόνο αν αλλάξουν τα ωμά δεδομένα -- παράγει
                            # ΚΑΙ τα δύο datasets (proportional, Shin)
sleeping_experts_experiment.py
calibration_analysis.py
calibration_correction.py
final_ranking.py
significance_test.py
fixed_share.py
contextual_experts.py
value_betting.py           # 6 πειράματα: σημειακός έλεγχος, πλήρης σάρωση,
                            # closing-on-closing, αιτιότητα panel, grid 108
                            # γραμμών, Shin έναντι basic value betting
value_betting_online_kelly.py
outcome_class_accuracy.py
temporal_trends.py
pooling_window_comparison.py
calibration_per_season.py
calibration_warmstart.py
market_devig_comparisons.py  # 4 πειράματα: opening/closing, συσχέτιση
                              # ποιότητας, Shin/basic, ECE/reliability
\end{verbatim}
Μόνο το \file{data\_processer.py} χρειάζεται ξανα-εκτέλεση αν αλλάξουν τα ωμά δεδομένα --
κάθε άλλο script ξανα-υπολογίζει τους δικούς του αριθμούς σε κάθε εκτέλεση.
(\file{sleeping\_experts.py} δεν εμφανίζεται εδώ -- είναι βιβλιοθήκη algorithms/data-loading
χωρίς δικό του \var{main()}, εισάγεται από όλα τα παραπάνω.)

\end{document}
